% claim-ref: eisenia-fetida#working-identity
% claim-ref: eisenia-fetida#feeding
% claim-ref: eisenia-fetida#moisture
% claim-ref: eisenia-fetida#routine-stop
% claim-ref: eisenia-fetida#hygiene-boundary
% binder-placement hero eisenia-fetida-pending; owner image pending; no substitute.
\begin{minipage}[t][3.30in][t]{3.36in}\vspace{0pt}\raggedright
\CompanionHeader{ANIMAL CARE}{Red wiggler}{Eisenia fetida / planning identity}{Also tiger worm. Surface-litter earthworm for a terrestrial bedding bin, not aquarium stocking. [EF-N2; EF-O1]}
{\sffamily\fontsize{9.5}{12}\selectfont Reference only: ownership, supplier, life stages and bin conditions unconfirmed.\par\vspace{6pt}Pale bands suggest E. fetida; E. andrei is usually uniformly red. Appearance is not genetic identification. Record supplier name and uncertainty. [EF-T2]\par\vspace{6pt}Owner photograph pending; no substitute.}
\end{minipage}\hfill\begin{minipage}[t][3.30in][t]{3.30in}\vspace{0pt}\AssetHero\end{minipage}\par\vspace{5pt}
\Context{Opaque, shaded bin; moist, breathable bedding. Gentle checks, not whole-bin turning. Published rates are not a ration for unmeasured stock. [EF-N1; EF-O1]}
\CardRow{Food and access}{Chopped fruit/veg, small coffee-ground portions, plain filters and finely crushed eggshells. Cover with about 5 cm (2 in) bedding. Check consumption before adding more; divert surplus. Bedding is also food. [EF-EPA; EF-N1]}{Conservative exclusions}{No meat, fish, bones, dairy, oils, pet waste, citrus, onion or garlic. No insecticide. For heavy insects, reduce excess feed and add bedding; small decomposers can be normal. [EF-EPA; EF-N1; EF-F1]}
\Panel{Moisture / trouble cues}{A squeeze should yield a few drops. Mist gradually if dry; add dry plain cardboard/newspaper and restore airflow/drainage if soggy. At feeding check leftovers, smell, temperature and escapes. Rotten odor: pause feed, remove unsuitable scraps and correct conditions. Persistent distress or losses: pause expansion and ask Extension. [EF-N1; proposed]}
\Panel{Hygiene and aquatic boundary}{Leachate may carry harmful microbes; it is not reliably safe worm tea. Keep off edible plant surfaces; wash hands. Keep bin tools separate. No leachate or compost in aquaria; this household protocol does not establish live-feed safety. [EF-N2; EF-F1; proposed]}
\Note{Care record}{Date / food remaining / moisture / odor / next check: \hrulefill}
\Sources{Clickable evidence: \href{https://www.epa.gov/recycle/composting-home}{EF-EPA}; \href{https://content.ces.ncsu.edu/worms-can-recycle-your-garbage}{EF-N1}; \href{https://content.ces.ncsu.edu/raising-earthworms-successfully}{EF-N2}; \href{https://extension.oregonstate.edu/catalog/em-9034-composting-worms}{EF-O1}; \href{https://blogs.ifas.ufl.edu/orangeco/2020/12/10/vermicomposting-using-worms-for-composting/}{EF-F1}; \href{https://journals.plos.org/plosone/article?id=10.1371/journal.pone.0191711}{EF-T2}. Claims, source scope, separate dates and contradictions: adjacent YAML. Planning only; habitat unmeasured.}{eisenia-fetida / animal-care}
